\subsubsection{Backend}
\subsubsubsection{Spring Boot}

\subsubsubsubsection{Giới thiệu}
Spring Boot là một framework Java mã nguồn mở, được phát triển từ hệ sinh thái Spring, thiết kế để đơn giản hóa quá trình phát triển ứng dụng thông qua cơ chế "Convention over Configuration"\ (Quy ước hơn cấu hình) \cite{walls-2022}. Nó cung cấp các tính năng tự động cấu hình (auto-configuration), tích hợp sẵn server (embedded server), và quản lý phụ thuộc tối ưu, giúp giảm thiểu các thiết lập rườm rà ban đầu \cite{spring-boot-ref}.
Trong ngữ cảnh của hệ thống gợi ý hành trình du lịch thông minh, Spring Boot đóng vai trò là xương sống (backbone) của phía server. Nó chịu trách nhiệm xử lý các luồng nghiệp vụ phức tạp bao gồm: tạo gợi ý hành trình dựa trên AI, quản lý định danh người dùng, tìm kiếm điểm đến và kiểm duyệt nội dung cộng đồng \cite{musib-2020}. Nền tảng này đảm bảo khả năng mở rộng (scalability) và tính toàn vẹn dữ liệu khi tích hợp với các nguồn thông tin bên ngoài \cite{webb-2020}.
\subsubsubsubsection{Ưu điểm}
\begin{itemize}
    \item Xử lý logic nghiệp vụ hiệu quả: 
    Spring Boot hỗ trợ phát triển các API REST nhanh chóng thông qua kiến trúc Spring MVC, giúp xử lý hiệu quả các yêu cầu phức tạp như gợi ý hành trình du lịch và cập nhật chi phí tức thời \cite{walls-2022}. Việc tách biệt rõ ràng giữa các lớp xử lý giúp mã nguồn dễ bảo trì và mở rộng.
    
    \item Tích hợp dữ liệu dễ dàng: 
    Các module chuyên biệt như Spring Batch và Spring Integration cung cấp khả năng thu thập và xử lý dữ liệu mạnh mẽ từ các nguồn bên ngoài (review, thời tiết, sự kiện...). Hệ thống hỗ trợ cả xử lý theo lô (batch) và theo luồng (streaming), đảm bảo dữ liệu luôn được chuẩn hóa và cập nhật ngữ cảnh thời gian thực \cite{spring-batch-ref}.
    
    \item Bảo mật mạnh mẽ: 
    Module Spring Security tích hợp sẵn giúp triển khai xác thực đa lớp (OAuth2/JWT) và phân quyền chi tiết. Điều này đóng vai trò then chốt trong việc bảo vệ dữ liệu nhạy cảm của người dùng như hồ sơ sở thích du lịch và lịch sử hội thoại \cite{spring-security-ref}.
    
    \item Quản lý người dùng và nội dung: 
    Spring Boot cung cấp nền tảng vững chắc để xây dựng các API quản trị, hỗ trợ các tác vụ kiểm duyệt nội dung cộng đồng, khóa tài khoản vi phạm và theo dõi vết (audit trail) hệ thống, đảm bảo tính minh bạch trong vận hành \cite{musib-2020}.
    
    \item Giám sát và bảo trì: 
    Công cụ Spring Boot Actuator được tích hợp sẵn cung cấp các chỉ số giám sát hiệu suất (metrics) và log sự kiện (health checks) mà không cần cấu hình phức tạp. Tính năng này giúp đội ngũ kỹ thuật dễ dàng theo dõi pipeline xử lý dữ liệu và hiệu quả của mô hình gợi ý trong thời gian thực \cite{webb-2020, spring-boot-ref}.
\end{itemize}

\subsubsubsubsection{So sánh Spring Boot và các công nghệ khác}
\clearpage
\begin{landscape}
\begin{table}[h!]
\centering
\setlength{\tabcolsep}{4pt}
\footnotesize
\begin{tabularx}{\linewidth}{|l|X|X|X|X|X|}
\hline
\textbf{Tiêu chí} & \textbf{Spring Boot} & \textbf{Node.js (Express)} & \textbf{Django (Python)} & \textbf{Ruby on Rails} & \textbf{Flask (Python)} \\
\hline
\textbf{Xử lý logic nghiệp vụ} & 
Cao: REST API dễ cấu hình, Spring MVC mạnh mẽ, hỗ trợ gợi ý hành trình và cập nhật chi phí tức thời \cite{walls-2022}. & 
Cao: Express nhẹ, nhanh \cite{express-docs}, phù hợp cho API gợi ý. Nhưng cần cài thêm thư viện cho logic phức tạp. & 
Cao: Django REST Framework mạnh mẽ \cite{django-drf}, hỗ trợ tốt việc quản lý nội dung. Cần cấu hình thêm cho AI. & 
Cao: Convention-over-configuration \cite{rails-guides}, phát triển nhanh nhưng chậm hơn Spring Boot ở các API phức tạp. & 
Trung bình: Vi khung (microframework) nhẹ, linh hoạt \cite{flask-docs}, nhưng cần tự xây dựng logic cấu trúc từ đầu. \\
\hline
\textbf{Tích hợp dữ liệu} & 
Cao: Spring Batch và Spring Integration hỗ trợ mạnh mẽ ETL (batch \& streaming) \cite{spring-batch-ref}. & 
Trung bình: Cần sử dụng các thư viện bên ngoài (như Axios, node-cron) để xử lý ETL, thiếu module chuẩn hóa tích hợp sẵn. & 
Cao: Hỗ trợ xử lý nền tốt thông qua Celery \cite{celery-docs} và Django ORM, nhưng streaming kém hơn Spring Integration. & 
Trung bình: Active Record xử lý dữ liệu tốt \cite{rails-guides}, nhưng streaming và ETL đòi hỏi cài đặt thêm các gem ngoài. & 
Trung bình: Yêu cầu cài đặt thêm các thư viện mở rộng như Flask-SQLAlchemy \cite{flask-sqlalchemy} để xử lý dữ liệu. \\
\hline
\textbf{Bảo mật} & 
Cao: Spring Security tích hợp sẵn OAuth2/JWT \cite{spring-security-ref}, bảo vệ tốt dữ liệu nhạy cảm. & 
Trung bình: Đòi hỏi cấu hình thủ công phức tạp qua các middleware bên ngoài như Passport.js \cite{passport-docs}. & 
Cao: Django Authentication có sẵn cơ chế bảo mật mạnh \cite{django-admin}, nhưng tích hợp OAuth2 tốn nhiều bước hơn. & 
Cao: Gem Devise \cite{devise-docs} cung cấp giải pháp xác thực toàn diện, nhưng tùy biến luồng OAuth2 phức tạp. & 
Trung bình: Cần phụ thuộc vào thư viện ngoài như Flask-Login \cite{flask-login}, không mạnh và đồng bộ như Spring Security. \\
\hline
\textbf{Quản trị người dùng/ nội dung} & 
Cao: Spring MVC hỗ trợ API quản trị, phân quyền và audit trail rõ ràng, minh bạch \cite{musib-2020}. & 
Trung bình: Không có sẵn module quản trị, lập trình viên phải tự thiết kế và xây dựng API từ đầu. & 
Cao: Tích hợp sẵn trang quản trị (Django Admin) \cite{django-admin} cực kỳ trực quan để kiểm duyệt nội dung. & 
Cao: Gem Rails Admin \cite{rails-admin} hỗ trợ tạo giao diện quản trị tự động rất nhanh chóng và tiện lợi. & 
Thấp: Hoàn toàn không có module quản trị tích hợp sẵn, đòi hỏi nỗ lực phát triển lớn từ con số không. \\
\hline
\textbf{Giám sát và bảo trì} & 
Cao: Spring Boot Actuator \cite{spring-boot-ref} cung cấp sẵn telemetry, log, health checks cực kỳ dễ theo dõi. & 
Trung bình: Cần cấu hình công cụ quản lý tiến trình bên ngoài như PM2 \cite{pm2-docs} và Winston để thu thập log. & 
Trung bình: Đòi hỏi tích hợp các hệ thống giám sát của bên thứ ba (như Sentry) để theo dõi chi tiết ứng dụng. & 
Trung bình: Cần sử dụng các công cụ hoặc gem ngoài (như New Relic) để có được chỉ số hiệu suất hệ thống. & 
Thấp: Cơ chế giám sát lỏng lẻo, lập trình viên phải tự thiết lập toàn bộ quy trình logging và monitoring. \\
\hline
\textbf{Hệ sinh thái và cộng đồng} & 
Cao: Hệ sinh thái Java khổng lồ, ổn định ở môi trường doanh nghiệp lớn \cite{spring-boot-ref}. & 
Cao: Hệ sinh thái npm phong phú nhưng phân tán, chất lượng thư viện không đồng đều \cite{express-docs}. & 
Cao: Cộng đồng Python lớn mạnh, hỗ trợ rất tốt nếu dự án có định hướng kết hợp AI/Machine Learning sâu. & 
Trung bình: Cộng đồng từng rất mạnh nhưng hiện đang thu hẹp hơn so với Node.js và Spring Boot. & 
Trung bình: Linh hoạt cho các dự án nhỏ hoặc microservices, nhưng thiếu tính đồng bộ cho dự án monolithic lớn. \\
\hline
\end{tabularx}
\caption{So sánh tính năng và độ tương thích của Spring Boot với các công nghệ Backend phổ biến}
\end{table}
\end{landscape}
